% TOP_PROBLEM: {"id": 2303007, "title": "Research Problems in Function Theory — Problem 3.7", "category_id": 8, "category": {"id": 8, "name": "analysis", "display_name": "Analysis", "description": "Limits, continuity, calculus, and function theory.", "slug": "analysis", "order_index": 8, "created_at": "2026-07-31T15:26:25.671Z"}, "status": "open"}
% FIRST_POSTED: null
% ENTRY_KIND: statement_only
% Imported from: batch12.tex; UnsolvedMath ID 2303007
% Compile twice: lualatex 2303007.tex
\documentclass[11pt]{article}
\usepackage{fontspec}
\setmainfont{Latin Modern Roman}
\usepackage{amsmath,amssymb,mathtools,unicode-math}
\setmathfont{Latin Modern Math}
\usepackage{xurl,hyperref}
\hypersetup{colorlinks=true,urlcolor=blue,pdftitle={UnsolvedMath Problem 2303007}}
\newcommand{\sref}[2]{\hyperlink{source:#1:#2}{[S#2]}}
\newcommand{\cataloguescope}[1]{\paragraph{Mathematical question.}#1}
\begin{document}
% UnsolvedMath ID 2303007; source catalogue code AMR-022-3007
\setcounter{section}{2303006}
\section{Research Problems in Function Theory — Problem 3.7}\label{2303007}
{\small\sffamily UnsolvedMath ID 2303007 · Analysis}
\subsection{Definitions and mathematical statement}

\paragraph{Shared notation.}$\mathbb N=\{1,2,\ldots\}$, and $\mathbb Z,\mathbb Q,\mathbb R,\mathbb C$ denote the integers, rationals, reals and complex numbers. Any other conventions are specified locally.


A subharmonic function is an upper semicontinuous function with values in $[-\infty,\infty)$, not identically $-\infty$ on a connected component, that satisfies the mean-value inequality on every closed ball in its domain. A real harmonic function is a twice continuously differentiable function $u$ satisfying $\Delta u=0$. For a subharmonic function $u$ on the plane set
\[B(r,u)=\sup_{|z|=r}u(z),\qquad T(r,u)=\frac1{2\pi}\int_0^{2\pi}\max\{u(re^{i\theta}),0\}\,d\theta.\]
Assume $u$ is nonconstant and its order $\rho=\limsup_{r\to\infty}\log T(r,u)/\log r$ is finite and nonnegative; $T(r,u)>0$ for all sufficiently large $r$. This is the subharmonic version of $\log M$ and the Nevanlinna characteristic in Problem 1.17.
The source update attributes the generalization to subharmonic functions in space to Dahlberg.
\paragraph{Statement recorded in the supplied source.}
Does the finite-order Paley bound extend to all such subharmonic functions: is there a finite constant $C(\rho)$ depending only on the order such that
\[1\leq\liminf_{r\to\infty}\frac{B(r,u)}{T(r,u)}\leq C(\rho)?\]
Determine the sharp constants. In the two ranges explicitly stated in Problem 1.17, does one have
\[C(\rho)=\frac{\pi\rho}{\sin(\pi\rho)}\quad(0<\rho<1/2),\qquad C(\rho)=\pi\rho\quad(\rho>1/2)?\]
No value at $\rho=0$ or $\rho=1/2$ is specified by these two displayed formulas. \sref{2303007}{2}
\subsection{Short English statement}
Does the sharp finite-order Paley inequality hold for arbitrary subharmonic functions, with the same constants as for logarithms of entire functions?
\subsection{Sources}
\begin{description}
\item[\hypertarget{source:2303007:1}{S1}] ULAM.AI and the UnsolvedMath Contributors, UnsolvedMath: A Curated Collection of Open Mathematics Problems, record 2303007 (AMR-022-3007). CC BY 4.0. \url{https://huggingface.co/datasets/ulamai/UnsolvedMath}.
\item[\hypertarget{source:2303007:2}{S2}] W. K. Hayman and E. F. Lingham, \emph{Research Problems in Function Theory}, arXiv:1809.07200v2 (2018), Problem 3.7 and Update 3.7. \url{https://arxiv.org/abs/1809.07200}.
\end{description}
\label{end:2303007}
\end{document}
